\subsection{Aplicaciones: I. Anillos japoneses locales.}
\label{IV.2.ss7.6-fr}

\begin{proposition}[7.6.1]
\label{IV.2.s7.6.1-fr}
Sea $A$ un anillo local noetheriano reducido cuyas fibras formales sean
geométricamente normales. Entonces el completado $\widehat A$ es reducido,
el cierre entero $A'$ de $A$ en su anillo total de fracciones es una
$A$-álgebra finita (y, por tanto, un anillo noetheriano semilocal), y su
completado $\widehat{A}'$ es isomorfo al cierre entero de $\widehat A$ en su
anillo total de fracciones.
\end{proposition}

\begin{proof}
Las fibras formales de $A$ son \emph{a fortiori} geométricamente reducidas;
por tanto, la hipótesis de que $A$ es reducido implica que también lo es
$\widehat A$, en virtud de (7.3.17). Sean $\mathfrak p_i$
($1\leq i\leq n$) los ideales primos minimales de $A$, y sea
$B_i=A/\mathfrak p_i$; las fibras formales de los anillos locales $B_i$
son entonces también geométricamente normales (7.3.15), por lo que los
$\widehat B_i$ son también reducidos. De (\textbf 0, 23.1.7, (i)) se deduce que
el cierre entero $B'_i$ de $B_i$ en su cuerpo de fracciones es un
$B_i$-módulo finitamente generado y, por tanto, un $A$-módulo finitamente
generado; como $A'$ es el producto directo de los $B'_i$ (\textbf{II}, 6.3.8),
se ve que $A'$ es un $A$-módulo finitamente generado. Sean
$\mathfrak m_j$ ($1\leq j\leq r$) los ideales maximales del anillo
semilocal $A'$; se sabe que el completado $\widehat{A}'$ de $A'$ se
identifica con el producto directo de los completados
$\widehat{A}'_{\mathfrak m_j}$ de los $A'_{\mathfrak m_j}$ (Bourbaki,
\oldpage[IV-2]{209}
\emph{Alg. comm.}, cap. III, § 2, n\textsuperscript{o} 13, cor. de la prop. 19).
Ahora bien, de la hipótesis y de (7.3.15) se deduce que las fibras formales
de los $A'_{\mathfrak m_j}$ son geométricamente normales; como
$\operatorname{Spec}(A')$ es normal por definición, también lo son los
$\operatorname{Spec}(A'_{\mathfrak m_j})$, y de (7.3.17) se deduce que
$\operatorname{Spec}(\widehat{A}'_{\mathfrak m_j})$ es normal para todo $j$,
y, por tanto, también lo es $\operatorname{Spec}(\widehat{A}')$.
Por otra parte (Bourbaki, \emph{Alg. comm.}, cap. IV, § 2,
n\textsuperscript{o} 5, cor. 3 de la prop. 9 y cap. III, § 3,
n\textsuperscript{o} 4, teorema 3), $\widehat{A}'$ se identifica con
$A'\otimes_A\widehat A$, puesto que $A'$ es un $A$-módulo finitamente
generado; como $A'$ contiene a $A$ y está contenido en el anillo total
de fracciones $R$ de $A$, y $\widehat A$ es un $A$-módulo plano,
$\widehat{A}'$ contiene a $\widehat A$ y está contenido en
$R'=R\otimes_A\widehat A$. Finalmente, como $\widehat A$ es un
$A$-módulo plano, todo elemento regular de $A$ es también
$\widehat A$-regular (\textbf 0$_{\mathrm I}$, 6.3.4); por tanto, $R'$ se
identifica canónicamente con un subanillo del anillo total de fracciones
$R''$ de $\widehat A$ (Bourbaki, \emph{Alg. comm.}, cap. II, § 2,
n\textsuperscript{o} 1, \emph{Observación} 7). Como
$\operatorname{Spec}(\widehat{A}')$ es normal y $\widehat{A}'$ es un
$\widehat A$-módulo finitamente generado, $\widehat{A}'$ es, en efecto,
el cierre entero de $\widehat A$ en $R''$.
\end{proof}
